\subsection{Spectrum-internal Features from Photon-Path Distribution Function}
\label{sec:methods_photon_path}

%\subsubsection{HEADING}

We derived spectral features based on the photon path distribution function (PPDF) assumptions that summarize the distribution of atmospheric
photon paths represented in each OCO-2 absorption band. The derivation follows
the radiative-transfer equivalence theorem
\citep{irvine1964,partain2000,stephens2000}. 


For a fixed illumination and viewing geometry with solar zenith angle (SZA) and viewing zenith angle (VZA), let $P_0(l)$ denote the ensemble of photon trajectories $l$ from the top of the atmosphere (TOA) reaching the detector in a conservative reference atmosphere with the same surface boundary and continuum scattering properties. The photon trajectories $l$ are mainly determined by multiple-scattering within the atmosphere by air molecules, clouds, and aerosols. Across each narrow OCO-2 band, we assume that these scattering properties vary slowly relative to molecular line absorption, so $P_0(l)$ is the same within each narrow band.  Wavelengths then changes the contribution $w_\lambda$ of a trajectory primarily through its Beer--Lambert survival probability,
%
\begin{equation}
  w_\lambda(l)
  = \exp[-\tau_\lambda(l)],
  \qquad
  \tau_\lambda(l)
  = \int_{l} \alpha_\lambda(s)\,\mathrm{d}s,
  \label{eq:main_path_weight}
\end{equation}
%
where $\alpha_\lambda$ is the local molecular absorption coefficient, and $\tau_\lambda$ is the total optical depth of each photon trajectory $l$.  The
normalized radiance is therefore the conservative-ensemble mean
%
\begin{equation}
  \mathcal T_\lambda
  = \int P_0(l)
      \exp[-\tau_\lambda(l)]\,\mathrm{d}l.
  \label{eq:main_equivalence}
\end{equation}
%
This construction separates the multiple-scattering calculation from molecular
absorption and remains applicable to three-dimensional cloud fields because the
cloud geometry is encoded in $P_0(l)$ \citep{stephens2000}.
%
%
%
For the direct geometrical slant path $L$, we define the modeled molecular slant optical depth $\mathrm{\overline{SOD}}_\lambda$ along $L$,
%
\begin{equation}
  \mathrm{\overline{SOD}}_\lambda
  = \int_0^{L} \tau_\lambda(s)\,\mathrm{d}s,
  \label{eq:main_SOD}
\end{equation}
%

To obtain a tractable scalar representation, we define the relative path length $l'$ for non-absorbing wavelength or an effective absorber-path enhancement relative to the direct geometrical slant as 

\begin{equation}
  l'
  = l/L,
  \label{eq:main_relative_path_length}
\end{equation}

Thus, we can approximate the absorber amount along each trajectory as
\begin{equation}
  A_\lambda(l')
  \simeq \mathrm{\overline{SOD}}_\lambda l',
  \label{eq:main_scalar_path}
\end{equation}

The conservative ensemble induces a normalized distribution $p_0(l')$, so Eq.~\eqref{eq:main_equivalence} reduces to

\begin{equation}
  \mathcal T(\mathrm{\overline{SOD}}_\lambda)
  = \int_0^\infty p_0(l')
      \exp(-\mathrm{\overline{SOD}}_\lambda l')\,\mathrm{d}l'.
  \label{eq:main_laplace}
\end{equation}

Thus, the resolved absorption spectrum samples the Laplace transform of one
wavelength-independent conservative path distribution within each band.  This
does not imply that weak and strong channels detect identical photon
populations: absorption reweights the detected paths as
$p_{\mathrm{det}}(l'\mid\mathrm{SOD})\propto
p_0(l')\exp(-\mathrm{SOD}\cdot l')$, preferentially suppressing long paths at large
$\mathrm{SOD}$.

We normalized the measured radiance as
\begin{equation}
  T_\lambda
  = \frac{\pi I_\lambda}{F_{0,\lambda}^{ILS}\mu_0}
  = \gamma\,\mathcal T(\tau_\lambda),
  \label{eq:main_transmittance}
\end{equation}
where $F_{0,\lambda}^{ILS}$ is the Earth-Sun distance corrected, Doppler stretch corrected, and
ILS-convolved top-of-atmosphere solar irradiance, $\mu_0$ is the
cosine of the SZA, and $\gamma$ is an effective continuum
reflectance.  In practice, $\gamma$ can include surface bidirectional
reflectance, polarization, and surface--atmosphere coupling and is not
interpreted as an exact Lambertian albedo.

Let $C_n$ denote the ordinary cumulants of $p_0(l')$.  Expanding the log
Laplace transform about zero absorption gives
\begin{equation}
  \ln T(\mathrm{\overline{SOD}})
  = \ln\gamma
    + \sum_{n=1}^{\infty}
      \frac{(-1)^n C_n}{n!}\tau^n.
  \label{eq:main_cumulant}
\end{equation}
For each sounding and band, we fitted the truncated form
\begin{equation}
  \ln T(\mathrm{\overline{SOD}})
  \simeq b
    + \sum_{n=1}^{N}
      \frac{(-1)^n k_n}{n}\tau^n,
  \qquad
  k_n=\frac{C_n}{(n-1)!},
  \label{eq:main_fit}
\end{equation}
where $b$ is the fitted continuum intercept.  Under the equivalence-theorem and
scalar-path assumptions, $k_1=\mathrm{E}_{p_0}[l']$ is the mean absorber-path
enhancement and $k_2=\mathrm{var}_{p_0}(l')$ is its variance; coefficients
$k_{n\geq3}$ are scaled higher cumulants.

The model is linear in $(b,k_1,\ldots,k_N)$.  We solved it by singular-value
least squares and used bounded-variable least squares only when the
unconstrained solution violated $k_1,k_2\geq0$.  We excluded the two edge
channels and channels with $T_\lambda>1$.  The truncation orders were 7, 3, and
7 for the O$_2$A, WCO$_2$, and SCO$_2$ bands, respectively; the lower
WCO$_2$ order reflects its smaller optical-depth range and the resulting
poor identifiability of high-order terms. 

We interpret these quantities as \emph{band-effective absorber-path cumulants}.
The scalar approximation in Eq.~\eqref{eq:main_scalar_path} can fail when
pressure- and temperature-dependent line absorption weights atmospheric layers
differently, or when surface, aerosol, or polarization properties vary within a
band.  A layer-resolved formulation and the distinction from a gamma path model
are given in Appendix~\ref{app:spectral_fit_derivation}.










% For a single photon with wavelength $\lambda$ traveling from the top of the atmosphere (TOA) to the satellite sensor along the path $L$, we can express its transmittance $T_\lambda(L)$ as

% \begin{equation}
%     T_\lambda(L)
%     =
%     \exp\left\{-\int_0^L\sum_i \left[\sigma_{i,\lambda}(l)n_i(l)\right]dl\right\}
%     =
%     \exp\left\{-\int_0^L\sum_i \left[\alpha_{i,\lambda}(l)\right]dl\right\}
%     =
%     \exp\left\{-\int_0^L\alpha_{t,\lambda}(l)dl\right\},
% \label{eq:single_photon_transmittance}
% \end{equation}

% where $l$ is the partial path along the total path $L$, $\sigma_{i,\lambda}(l)$ is the extinction coefficient at wavelength $\lambda$ on the path $l$ for gas $i$, $n_i(l)$ is the number concentration of the path $l$ for gas $i$. $\alpha_{i,\lambda}$ is the extinction at $\lambda$ for gas $i$, while $\alpha_{t,\lambda}$ is the total extinction considering all gases.


% For multiple photons, we can calculate an average transmittance $\overline{T_\lambda(\alpha_{t,p}, L_p)}$ as 



% \begin{equation}
% \begin{split}
%     \overline{T_\lambda(\alpha_{t,p}, L_p)}
%     &=
%     \frac{1}{N}\sum_{p=1}^{N}
%     \exp\left\{-\int_0^{L_p}\alpha_{t,p}(l)dl\right\} \\
%     &\approx
%     \frac{1}{N}\sum_{p=1}^{N}
%     \exp\left\{-\int_0^{L_p}\overline{\alpha_t}dl\right\}
%     &=
%     \frac{1}{N}\sum_{p=1}^{N}
%     \exp\left\{\overline{\alpha_t}L_p\right\}
% \end{split}
% \label{eq:multi_photon_transmittance}
% \end{equation}


% where $\alpha_{t,p}(l)$ is the total extinction of the photon $p$ at the partial path $l$, and $L_p$ is the total path length of the photon $p$, $N$ is the total photon number. $\overline{\alpha_t}$ is the vertical mean total extinction, which is independent on photons. \\


% Equation \ref{eq:multi_photon_transmittance} can be expressed in the photon path distribution function $p(L)$


% \begin{equation}
% \begin{split}
%     \overline{T_\lambda(\alpha_{t,p}, L_p)}
%     &=
%     \int_0^{L\to \infty} p(L)
%     \exp\left\{-\overline{\alpha_t}L\right\}dL,
% \end{split}
% \label{eq:multi_photon_transmittance_2}
% \end{equation}

% Now we can define the relative photon path length $L'=L/L_{geo}$, where $L_{geo}$ is the reference geometric photon path length, or the vertical distance from TOA to the surface multiply by the geometric two-way airmass $m$,

% \begin{equation}
% m
% =
% \frac{1}{\cos\theta} + \frac{1}{\cos\phi} ,
% \label{eq:airmass_sod}
% \end{equation}

% where $\theta$ and $\phi$ are the solar and viewing zenith angles.

% \begin{equation}
% L_{geo}
% =
% m\,L_{\mathrm{TOA-surface}} ,
% \label{eq:airmass_photon_path_geo}
% \end{equation}


% \begin{equation}
% \begin{split}
%     \overline{T_\lambda(\alpha_{t,p}, L_p)}
%     &=
%     \int_0^\infty p(L)
%     \exp\left\{-\overline{\alpha_t}L\right\}dL,
% \end{split}
% \label{eq:multi_photon_transmittance_2}
% \end{equation}


% \begin{equation}
% \overline{\mathrm{SOD}}_c(\lambda)
% =
% \bar{\alpha}(\lambda)\bar{L},
% \label{eq:slant_optical_depth}
% \end{equation}


% For each footprint and spectral band, we use an
% instrument lineshape-convolved effective slant optical depth,
% $\overline{\mathrm{SOD}}_c$, as the spectral coordinate. The high-resolution
% vertical optical depth ($\tau_v(\lambda)\}$) is first multiplied by the geometric two-way airmass,






% The slant optical depth is then computed in transmission space,
% \begin{equation}
% \overline{\mathrm{SOD}}_j
% =
% -\ln
% \left[
% \frac{
% \sum_q R_j(\lambda_q)
% \exp\{-m\tau_v(\lambda_q)\}
% }{
% \sum_q R_j(\lambda_q)
% }
% \right],
% \label{eq:ils_sod}
% \end{equation}
% where $j$ indexes the discrete wavelength samples in an OCO-2 band and
% $q$ indexes the high-resolution wavelength grid used for the ILS convolution. $R_j(\lambda)$ is the instrument lineshape (ILS) for pixel or discrete wavelength $j$ and $\tau_v(\lambda_q)$ is the high-resolution vertical optical depth, including molecular absorption and Rayleigh extinction.




% The measured spectral transmittance used in the fit is
% \begin{equation}
% T(\lambda)
% =
% \frac{\pi I(\lambda)}{\mathcal{F}_{0}^{ILS}(\lambda)\cos\theta},
% \label{eq:measured_transmittance}
% \end{equation}

% where $I(\lambda)$ is the calibrated OCO-2 radiance and $\mathcal{F}_{0}^{ILS}(\lambda)$ is the Earth-Sun distance corrected, Doppler stretch corrected, and
% ILS-convolved solar normalization used in the spectral fit. See details in \ref{eq:solar_distance_doppler}, \ref{eq:solar_ils_convolution}, \ref{eq:effective_solar_normalization}.











% \begin{equation}
% T(\overline{\mathrm{SOD}}_c)
% =
% \gamma
% \int_0^\infty
% p(l')
% \exp(-\overline{\mathrm{SOD}}_c \; l')
% \,\mathrm{d}l' ,
% \label{eq:transmittance_laplace}
% \end{equation}

% \begin{equation}
% \ln T(\overline{\mathrm{SOD}}_c)
% =
% \ln\gamma
% +
% \sum_{n=1}^{N}
% (-1)^n
% \frac{k_n}{n}
% \overline{\mathrm{SOD}}_c^n .
% \label{eq:cumulant_fit}
% \end{equation}


% \begin{equation}
% k_1 = \langle l' \rangle,
% \qquad
% k_2 = \operatorname{var}(l'),
% \qquad
% \kappa = \frac{k_1^2}{k_2}.
% \label{eq:coefficient_interpretation}
% \end{equation}